% GENERATED FILE. Edit canonical lessons, then run npm run generate:books.
% canonical-source-hash: fnv1a64:bd53ca1b0a067185
% canonical-lessons: LA-C306-caedes, LA-C306-homicida, LA-C306-sicarius, LA-C306-pugio, LA-C306-occido, LA-R306-review-299-306

\chapter{Caedēs}
\label{ch:la-caedes}
\hlchaptermodality{306}

\begin{chapteropening}
\textbf{By the end of this chapter:} \emph{I can talk about violent crime: slaughter, a murderer, an assassin, a dagger and to kill.}

Complete the last lesson of chapter 306: I can talk about violent crime: slaughter, a murderer, an assassin, a dagger and to kill.
\end{chapteropening}

\section[caedēs, caedis]{caedēs, caedis — slaughter, murder}

\label{lesson:LA-C306-caedes}



\begin{quote}
Before the new one: say the Latin for fraud, then the Latin for a trick.
\end{quote}

\subsection*{You'll want to know: caedēs, caedis}
\textbf{caedēs, caedis} — \textquotedblleft{}slaughter, murder\textquotedblright{}.

\begin{quote}
Post caedem tōta urbs tacita erat.
\end{quote}

\begin{quote}
\emph{After the murder the whole city was silent.}
\end{quote}

Feminine, 3rd declension, genitive plural \emph{caedium}. Used of one killing or of a massacre.

\begin{cousinweb}
\textbf{caedēs} is built on \emph{caedō}, \textquotedblleft{}to cut, to cut down\textquotedblright{} (EtymDB), which is inherited from a root meaning \textquotedblleft{}to cut, to hew\textquotedblright{}.
\end{cousinweb}

\begin{grammarlens}[title={What you've built}]
One more word for telling, explaining and describing at B1.
\end{grammarlens}

\subsection*{Your turn}
\begin{itemize}
\raggedright
  \item \emph{Say it:} \emph{caedēs, caedis}
  \item \emph{Say it:} \emph{caedēs, caedis}, once more
  \item \emph{Say it:} a sentence of your own with \emph{caedēs, caedis}
  \item \emph{Recall:} say the Latin for fraud, then the Latin for a trick, then say \emph{caedēs, caedis} again
\end{itemize}

\begin{tcolorbox}[breakable,colback=teal!4,colframe=teal!35!black,title={Before you move on}]
What does caedēs, caedis mean? (\textquotedblleft{}Slaughter, murder\textquotedblright{}.) Say it once more.
\end{tcolorbox}
\section[homicīda, homicīdae]{homicīda, homicīdae — a murderer, a killer}

\label{lesson:LA-C306-homicida}



\begin{quote}
Before the new one: say the Latin for a trick, then the Latin for slaughter.
\end{quote}

\subsection*{You'll want to know: homicīda, homicīdae}
\textbf{homicīda, homicīdae} — \textquotedblleft{}a murderer, a killer\textquotedblright{}.

\begin{quote}
Homicīda ē carcere fūgit.
\end{quote}

\begin{quote}
\emph{The murderer escaped from prison.}
\end{quote}

Masculine or feminine, 1st declension. English \emph{homicide} names both the deed and the doer.

\begin{cousinweb}
\textbf{homicīda} joins \emph{homō}, \textquotedblleft{}person\textquotedblright{}, and \emph{-cīda}, \textquotedblleft{}killer\textquotedblright{}, from \emph{caedō}, \textquotedblleft{}to cut down\textquotedblright{} (Lewis \& Short). English \emph{homicide} comes from it.
\end{cousinweb}

\begin{grammarlens}[title={What you've built}]
One more word for telling, explaining and describing at B1.
\end{grammarlens}

\subsection*{Your turn}
\begin{itemize}
\raggedright
  \item \emph{Say it:} \emph{homicīda, homicīdae}
  \item \emph{Say it:} \emph{homicīda, homicīdae}, once more
  \item \emph{Say it:} a sentence of your own with \emph{homicīda, homicīdae}
  \item \emph{Recall:} say the Latin for a trick, then the Latin for slaughter, then say \emph{homicīda, homicīdae} again
\end{itemize}

\begin{tcolorbox}[breakable,colback=teal!4,colframe=teal!35!black,title={Before you move on}]
What does homicīda, homicīdae mean? (\textquotedblleft{}A murderer, a killer\textquotedblright{}.) Say it once more.
\end{tcolorbox}
\section[sīcārius, sīcāriī]{sīcārius, sīcāriī — an assassin}

\label{lesson:LA-C306-sicarius}



\begin{quote}
Before the new one: say the Latin for slaughter, then the Latin for a murderer.
\end{quote}

\subsection*{You'll want to know: sīcārius, sīcāriī}
\textbf{sīcārius, sīcāriī} — \textquotedblleft{}an assassin\textquotedblright{}.

\begin{quote}
Sīcāriī cōnsulem in viā exspectābant.
\end{quote}

\begin{quote}
\emph{The assassins were waiting for the consul in the street.}
\end{quote}

Masculine. Murder trials at Rome were held \emph{inter sīcāriōs}: among the assassins.

\begin{cousinweb}
\textbf{sīcārius} is built on \emph{sīca}, \textquotedblleft{}a curved dagger\textquotedblright{} (Lewis \& Short): the dagger man.
\end{cousinweb}

\begin{grammarlens}[title={What you've built}]
One more word for telling, explaining and describing at B1.
\end{grammarlens}

\subsection*{Your turn}
\begin{itemize}
\raggedright
  \item \emph{Say it:} \emph{sīcārius, sīcāriī}
  \item \emph{Say it:} \emph{sīcārius, sīcāriī}, once more
  \item \emph{Say it:} a sentence of your own with \emph{sīcārius, sīcāriī}
  \item \emph{Recall:} say the Latin for slaughter, then the Latin for a murderer, then say \emph{sīcārius, sīcāriī} again
\end{itemize}

\begin{tcolorbox}[breakable,colback=teal!4,colframe=teal!35!black,title={Before you move on}]
What does sīcārius, sīcāriī mean? (\textquotedblleft{}An assassin\textquotedblright{}.) Say it once more.
\end{tcolorbox}
\section[pugiō, pugiōnis]{pugiō, pugiōnis — a dagger}

\label{lesson:LA-C306-pugio}



\begin{quote}
Before the new one: say the Latin for a murderer, then the Latin for an assassin.
\end{quote}

\subsection*{You'll want to know: pugiō, pugiōnis}
\textbf{pugiō, pugiōnis} — \textquotedblleft{}a dagger\textquotedblright{}.

\begin{quote}
Pugiōnem sub veste habēbat.
\end{quote}

\begin{quote}
\emph{He had a dagger under his clothing.}
\end{quote}

Masculine, 3rd declension. A short stabbing blade that soldiers wore on the belt.

\begin{cousinweb}
\textbf{pugiō} is built, says Lewis \& Short, on the root of \emph{pungō}, \textquotedblleft{}to prick, to stab\textquotedblright{}; an old explanation says the dagger strikes with the point.
\end{cousinweb}

\begin{grammarlens}[title={What you've built}]
One more word for telling, explaining and describing at B1.
\end{grammarlens}

\subsection*{Your turn}
\begin{itemize}
\raggedright
  \item \emph{Say it:} \emph{pugiō, pugiōnis}
  \item \emph{Say it:} \emph{pugiō, pugiōnis}, once more
  \item \emph{Say it:} a sentence of your own with \emph{pugiō, pugiōnis}
  \item \emph{Recall:} say the Latin for a murderer, then the Latin for an assassin, then say \emph{pugiō, pugiōnis} again
\end{itemize}

\begin{tcolorbox}[breakable,colback=teal!4,colframe=teal!35!black,title={Before you move on}]
What does pugiō, pugiōnis mean? (\textquotedblleft{}A dagger\textquotedblright{}.) Say it once more.
\end{tcolorbox}
\section[occīdō, occīdere]{occīdō, occīdere — to kill, to cut down}

\label{lesson:LA-C306-occido}



\begin{quote}
Before the new one: say the Latin for an assassin, then the Latin for a dagger.
\end{quote}

\subsection*{You'll want to know: occīdō, occīdere}
\textbf{occīdō, occīdere} — \textquotedblleft{}to kill, to cut down\textquotedblright{}.

\begin{quote}
Latrō viātōrem occīdit et pecūniam rapuit.
\end{quote}

\begin{quote}
\emph{The bandit killed a traveller and snatched his money.}
\end{quote}

\emph{occīdō, occīdere, occīdī, occīsum}, with a long \emph{ī}. Do not confuse \emph{occidō}, 'to fall, to set', of the sun.

\begin{cousinweb}
\textbf{occīdō} is built on \emph{caedō}, \textquotedblleft{}to cut down\textquotedblright{}, with \emph{ob-} (EtymDB, Lewis \& Short): to strike down. Its look-alike \emph{occidō}, \textquotedblleft{}to fall\textquotedblright{}, is built on \emph{cadō}.
\end{cousinweb}

\begin{grammarlens}[title={What you've built}]
One more word for telling, explaining and describing at B1.
\end{grammarlens}

\subsection*{Your turn}
\begin{itemize}
\raggedright
  \item \emph{Say it:} \emph{occīdō, occīdere}
  \item \emph{Say it:} \emph{occīdō, occīdere}, once more
  \item \emph{Say it:} a sentence of your own with \emph{occīdō, occīdere}
  \item \emph{Recall:} say the Latin for an assassin, then the Latin for a dagger, then say \emph{occīdō, occīdere} again
\end{itemize}

\begin{tcolorbox}[breakable,colback=teal!4,colframe=teal!35!black,title={Before you move on}]
What does occīdō, occīdere mean? (\textquotedblleft{}To kill, to cut down\textquotedblright{}.) Say it once more.
\end{tcolorbox}
\section[(dialogue)]{Words from chapters 299–306 — a review}

\label{lesson:LA-R306-review-299-306}



\begin{quote}
Say the Latin for a narrative and to recall.
\end{quote}

\subsection*{Your turn}
Say these aloud:

\begin{itemize}
\raggedright
  \item a narrative, to recall, to tell of, to report, I say
  \item a letter-carrier, a herald, an instruction, an embassy, to bring back word
  \item a whisper, to whisper, to spread abroad, to grow stronger, to assert
  \item a spy, a scout, to search out, to betray, to open up
  \item a theft, highway robbery, a mugger, a robber, loot
  \item to break open, a bolt, to snatch, to catch, to seize
  \item a crime, a deed, wicked, fraud, a trick
  \item slaughter, a murderer, an assassin, a dagger, to kill
\end{itemize}

\begin{tcolorbox}[breakable,colback=teal!4,colframe=teal!35!black,title={Before you move on}]
A narrative? (\textbf{nārrātiō, nārrātiōnis}.) To recall? (\textbf{commemorō, commemorāre}.) To tell of? (\textbf{memorō, memorāre}.) To report? (\textbf{referō, referre}.) I say? (\textbf{inquam}.) A letter-carrier? (\textbf{tabellārius, tabellāriī}.) A herald? (\textbf{praecō, praecōnis}.) An instruction? (\textbf{mandātum, mandātī}.) An embassy? (\textbf{lēgātiō, lēgātiōnis}.) To bring back word? (\textbf{renūntiō, renūntiāre}.) A whisper? (\textbf{susurrus, susurrī}.) To whisper? (\textbf{susurrō, susurrāre}.) To spread abroad? (\textbf{vulgō, vulgāre}.) To grow stronger? (\textbf{crēbrēscō, crēbrēscere}.) To assert? (\textbf{perhibeō, perhibēre}.) A spy? (\textbf{speculātor, speculātōris}.) A scout? (\textbf{explōrātor, explōrātōris}.) To search out? (\textbf{explōrō, explōrāre}.) To betray? (\textbf{prōdō, prōdere}.) To open up? (\textbf{patefaciō, patefacere}.) A theft? (\textbf{fūrtum, fūrtī}.) Highway robbery? (\textbf{latrōcinium, latrōciniī}.) A mugger? (\textbf{grassātor, grassātōris}.) A robber? (\textbf{praedō, praedōnis}.) Loot? (\textbf{praeda, praedae}.) To break open? (\textbf{effringō, effringere}.) A bolt? (\textbf{claustrum, claustrī}.) To snatch? (\textbf{rapiō, rapere}.) To catch? (\textbf{dēprehendō, dēprehendere}.) To seize? (\textbf{comprehendō, comprehendere}.) A crime? (\textbf{scelus, sceleris}.) A deed? (\textbf{facinus, facinoris}.) Wicked? (\textbf{scelerātus, scelerāta, scelerātum}.) Fraud? (\textbf{fraus, fraudis}.) A trick? (\textbf{dolus, dolī}.) Slaughter? (\textbf{caedēs, caedis}.) A murderer? (\textbf{homicīda, homicīdae}.) An assassin? (\textbf{sīcārius, sīcāriī}.) A dagger? (\textbf{pugiō, pugiōnis}.) To kill? (\textbf{occīdō, occīdere}.)
\end{tcolorbox}
